% Draft: environmental impact of Strategy A and Strategy B (v5, airports, context on).
% Numbers from notebooks/v5_eis.ipynb (built by notebooks/build_v5_eis_notebook.py); figures in notebooks/figures/v5_eis/.
% Citation keys to check against the paper's .bib: bahi2026environmental (already cited), ecologits (the EcoLogits
% reference of the DOLAP study), google2009search (new, entry at the end of this file).

%% ------------------------------------------------------------------------------------------------ Metrics
\subsubsection{Environmental impact}

\paragraph{Energy of a run.}
Both strategies are measured with the same formula, the agentic energy multiplier of Bahi et
al.~\cite{bahi2026environmental}. The energy of a run is
\[
E_{\text{run}} = \left(N_{\text{calls}} + N_{\text{tools}} \cdot w_{\text{tool}}\right) \cdot E_{\text{single}},
\]
where:
\begin{itemize}
\item $N_{\text{calls}}$ is the number of LLM invocations of the run;
\item $N_{\text{tools}}$ is the number of tool calls (here, web searches);
\item $E_{\text{single}}$ is the energy of one LLM invocation;
\item $w_{\text{tool}}$ is the energy of one tool call relative to one LLM invocation.
\end{itemize}
$E_{\text{single}}$ is the mean EcoLogits~\cite{ecologits} energy of one LLM call of the given model and strategy, so
that $N_{\text{calls}} \cdot E_{\text{single}}$ is the EcoLogits energy of the run's calls.

Bahi et al. give $w_{\text{tool}}$ only qualitatively, placing search over very large corpora among the most
energy-intensive tool calls. We derive it from the energy of one web search, $e_{\text{search}}$:
\[
w_{\text{tool}} = \frac{e_{\text{search}}}{E_{\text{single}}}
\qquad\Longrightarrow\qquad
E_{\text{run}} = N_{\text{calls}} \cdot E_{\text{single}} + N_{\text{tools}} \cdot e_{\text{search}} .
\]
The search provider does not report the energy of its searches. We use $e_{\text{search}} = 0.3$\,Wh, the only public
per-query figure for a web search engine~\cite{google2009search}. As this figure is dated and refers to another
service, we also report results for $e_{\text{search}}$ between 0 and 1\,Wh. Searches are counted whether or not our
local search cache answered them, since a user without the cache would make them.

The two strategies differ only in the values the formula takes:
\begin{itemize}
\item \textbf{Strategy~A} makes one call per run (1.00--1.08 on average, with the repair calls of unparsable answers) and
no tool calls. Its $E_{\text{run}}$ is therefore the EcoLogits energy of its calls.
\item \textbf{Strategy~B} makes two calls (planning and synthesis) and 1.2--3.0 searches per run on average. Its
$E_{\text{single}}$ is larger than in Strategy~A, because the synthesis call reads the retrieved evidence.
\end{itemize}

\paragraph{Composite environmental impact score.}
As in the DOLAP study, we compute three indicators for each configuration (model and prompt for Strategy~A, model and
variant for Strategy~B): energy consumption (EC, Wh), global warming potential (GWP, g\,CO$_2$e) and water consumption
footprint (WCF, mL), each the mean over the configuration's runs. All three use the same formula as the energy:
\[
I_{\text{run}} = N_{\text{calls}} \cdot I_{\text{single}} + N_{\text{tools}} \cdot i_{\text{search}},
\qquad I \in \{\mathrm{EC},\,\mathrm{GWP},\,\mathrm{WCF}\},
\]
where $I_{\text{single}}$ is the EcoLogits estimate of the indicator for one LLM call and $i_{\text{search}}$ is its
value for one search. For EC, $i_{\text{search}} = e_{\text{search}}$. A search has no GWP or WCF estimate, so we take
the GWP and WCF per Wh of the run's own LLM calls, i.e.\ we assume that a search runs in a data centre comparable to the
model's. Each indicator $I$ is normalised with a log-based min--max over all configurations of \emph{both} strategies,
and the three are averaged without weights:
\[
\mathrm{EIS} = \frac{1}{3} \sum_{I \in \{\mathrm{EC},\,\mathrm{GWP},\,\mathrm{WCF}\}}
\frac{\ln I - \ln I_{\min}}{\ln I_{\max} - \ln I_{\min}} .
\]
Normalising both strategies together places them on a common scale: 0 is the configuration with the lowest impact in
this study and 1 the one with the highest. The score is therefore relative to the configurations compared and cannot be
compared number for number with the scores of the DOLAP study, which were normalised over different models; absolute
values (Wh per run and per correct answer) are reported alongside it.

%% ------------------------------------------------------------------------------------------------ Results
\subsection{Environmental impact}

Table~\ref{tab:eis} compares, for each model, its best Strategy~A prompt with the agent, and
Figure~\ref{fig:eis-energy} shows the energy per run of every approach.

\begin{table}[t]
\centering
\label{tab:eis}
\small
\begin{tabular}{lrrrrrrrrr}
\toprule
& \multicolumn{2}{c}{EC per run (Wh)} & & Searches' & \multicolumn{2}{c}{EC per correct (Wh)} & \multicolumn{2}{c}{EIS} & \\
\cmidrule(lr){2-3}\cmidrule(lr){6-7}\cmidrule(lr){8-9}
Model & A & B & B/A & share & A & B & A & B & F1: A / B \\
\midrule
gpt-oss:20b         & 0.009  & 0.69 & 75$\times$  & 97\% & 0.22 & 10.5 & 0.33 & 0.89 & 0.161 / 0.249 (0.38 / 0.50) \\
gemma4:31b          & 0.0007 & 0.36 & 521$\times$ & 98\% & 0.02 & 10.5 & 0.00 & 0.81 & 0.190 / 0.152 (0.48 / 0.46) \\
gpt-oss:120b        & 0.033  & 0.97 & 29$\times$  & 91\% & 0.44 & 11.9 & 0.49 & 0.93 & 0.237 / 0.316 (0.40 / 0.62) \\
mistral-large-3     & 0.13   & 1.67 & 13$\times$  & 54\% & 1.96 & 20.7 & 0.62 & 0.94 & 0.219 / 0.294 (0.39 / 0.54) \\
deepseek-v4.1-flash & 0.047  & 1.27 & 27$\times$  & 68\% & 0.69 & 34.4 & 0.49 & 0.89 & 0.231 / 0.163 (0.43 / 0.47) \\
\bottomrule
\end{tabular}
\end{table}

\begin{figure}[t]
\centering
\includegraphics[width=\linewidth]{figures/v5_eis/eis_energy_per_run.pdf}
\caption{Energy per run (log scale) of each model with Strategy~A (zero-shot and best prompt) and Strategy~B (search,
search and memory). The open diamond is the agent's LLM energy alone; the bars give the share of the agent's energy
due to its web searches ($e_{\text{search}} = 0.3$\,Wh).}
\label{fig:eis-energy}
\end{figure}

\paragraph{The agent has a much higher impact.}
Per run, the agent uses 0.36--1.67\,Wh, against 0.0007--0.13\,Wh for the best Strategy~A prompt of the same model:
13 (mistral-large-3) to 521 (gemma4:31b) times more. Its LLM calls alone account for a factor of 2.3--11, because the
agent makes two calls instead of one and its synthesis call reads about 3{,}500 input tokens of evidence, against about
400 for Strategy~A. The rest comes from the searches, which make up 91--98\% of the agent's energy for the two gpt-oss
models and gemma4:31b, and 54--68\% for mistral-large-3 and deepseek-v4.1-flash: a search costs the same whichever
model requests it, so it weighs most on the agents built on small models. Per correct answer, the agent needs
10--34\,Wh, against 0.02--2.0\,Wh for the best Strategy~A prompt. Allowing the agent to use its own knowledge adds no
search and almost no energy, and it lowers the energy per correct answer of the models whose accuracy it raises
(gemma4:31b from 10.6 to 6.9\,Wh, deepseek-v4.1-flash from 34.4 to 21.6\,Wh).

\paragraph{Composite score.}
Over the 55 configurations, Strategy~A scores 0.00--0.62 and Strategy~B 0.81--0.94 (Figure~\ref{fig:eis-composite}):
every agent configuration has a higher impact than every LLM-only configuration. Within each strategy the model
remains the main factor, from gemma4:31b (lowest) to mistral-large-3 (highest). GWP does not follow EC exactly across
models: EcoLogits assigns mistral-large-3 and deepseek-v4.1-flash a lower-carbon electricity mix, so that their GWP per
Wh is a fifth to a quarter of that of the gpt-oss models. Counting the LLM calls alone, the two strategies would
overlap (Strategy~B: 0.36--1.00): the separation rests on the energy of the searches.

\begin{figure}[t]
\centering
\includegraphics[width=\linewidth]{figures/v5_eis/eis_composite_heatmap.pdf}
\caption{Composite EIS of each model and approach, normalised over the 55 configurations of both strategies, with the
web searches (left) and with the LLM calls only (right).}
\label{fig:eis-composite}
\end{figure}

\paragraph{Impact against accuracy.}
Figure~\ref{fig:eis-f1} relates the EIS of every configuration to its accuracy. For gpt-oss:20b, gpt-oss:120b and
mistral-large-3, the agent reaches the highest F1 (0.25--0.32, against 0.16--0.24 for their best Strategy~A prompt) at
the highest impact. For gemma4:31b and deepseek-v4.1-flash, the agent with search alone has both a higher impact and a
lower F1 than the best Strategy~A prompt; allowing memory brings their F1 level with Strategy~A (0.21 and 0.24), at the
agent's impact. With abstentions ignored, the agent is ahead for four models (0.50--0.62, and 0.53 for
deepseek-v4.1-flash with memory, against 0.38--0.43) and level for gemma4:31b (0.46 against 0.48). The configuration
with the lowest impact, gemma4:31b with Strategy~A, reaches F1 0.19 (0.48 with abstentions ignored), more than
gpt-oss:20b with any Strategy~A prompt.

\begin{figure}[t]
\centering
\includegraphics[width=\linewidth]{figures/v5_eis/eis_vs_f1.pdf}
\caption{F1 (left) and F1 with abstentions ignored (right) against the composite EIS, for the 45 Strategy~A
configurations (small markers; squares: each model's best prompt) and the 10 Strategy~B configurations (diamonds:
search; triangles: search and memory).}
\label{fig:eis-f1}
\end{figure}

\paragraph{Sensitivity to the energy of a search.}
The size of the gap depends on $e_{\text{search}}$, which we cannot measure (Figure~\ref{fig:eis-sensitivity}). With
the LLM calls only ($e_{\text{search}} = 0$), the agent uses 2--11 times the energy of the best Strategy~A prompt; at
$e_{\text{search}} = 0.03$\,Wh, ten times below the Google figure, 5--62 times; at 1\,Wh, 29--1{,}710 times. The
searches outweigh the LLM calls from 6--30\,mWh per search for the agents built on gemma4:31b and the gpt-oss models,
and from 0.14--0.26\,Wh for deepseek-v4.1-flash and mistral-large-3. The conclusion that the agent has a higher impact
than the LLM alone therefore holds for every model and every value of $e_{\text{search}}$; its magnitude does not, and
a direct measurement of the energy of web search, e.g.\ on a self-hosted search engine, is needed to fix it.

\begin{figure}[t]
\centering
\includegraphics[width=\linewidth]{figures/v5_eis/eis_search_sensitivity.pdf}
\caption{Energy per run of the agent (search) relative to the best Strategy~A prompt (left) and share of the agent's
energy due to web searches (right), as functions of the energy of one search. Dashed line: $e_{\text{search}} = 0.3$\,Wh.}
\label{fig:eis-sensitivity}
\end{figure}

%% ------------------------------------------------------------------------------------------------ bib entry (new)
% @misc{google2009search,
%   author       = {{Google}},
%   title        = {Powering a {Google} search},
%   howpublished = {Official Google Blog},
%   year         = {2009},
%   month        = jan,
%   note         = {11 January 2009. \url{https://googleblog.blogspot.com/2009/01/powering-google-search.html}}
% }
